\section{Caesarea: Workshop, Classroom, and Service to Churches}
\label{sec:origen-caesarea}

\subsection{A new center beside Jewish and Christian libraries}

Caesarea in Palestine became Origen's principal home from the early 230s. It
placed him among Alexander of Jerusalem, Theoctistus of Caesarea, Firmilian of
Cappadocia, Ambrose, clergy, visiting students, Jewish interlocutors, and
manuscript collections. The library later developed by Pamphilus and used by
Eusebius made Origen's papers and books part of an institutional memory. That
later continuity is one reason Eusebius knows so much and also why his life is
a defense of the Caesarean inheritance.

The comparative biblical project did not begin with this move. Eusebius places
his account of Hebrew and the Greek versions associated with Aquila,
Symmachus, Theodotion, and others before the permanent departure
(\emph{Ecclesiastical History} 6.16--17). He also says that a woman named
Juliana supplied works of Symmachus inherited from the translator, but gives no
secure location for the transfer. Modern reconstructions reasonably place
continued development of the project amid the books and Jewish-Christian
contacts available in Palestine; the surviving evidence does not assign each
acquisition or column to Alexandria or Caesarea.

Two exact texts control the ethical limit used here. In \emph{Letter to
Africanus} 4--9 Origen reports Hebrew consultation and extensive comparison
while also alleging that Jewish authorities removed passages discreditable to
their elders. In \emph{Against Celsus} 2.4--8 he calls Jewish Scripture the
foundation honored by Christianity, then describes contemporary Jews as
abandoned and their displacement as punishment for rejecting Jesus. These
passages establish learned dependence, supersession, and collective blame in
the checked argument. This publication does not convert them into a complete
lexicon or claim that every reference to Jews has the same speaker, audience,
or force.

\subsection{Gregory's portrait of a teacher}

The strongest near-contemporary portrait comes from a pupil traditionally
identified as Gregory Thaumaturgus. His farewell address says that he and his
brother had intended to study Roman law at Berytus but remained with Origen for
about five years. Panegyric magnifies its subject, yet its curricular detail is
biographically rare.

Origen first attended to a student's character, using questions and argument
to expose untested assumptions. He led pupils through logic, natural inquiry,
geometry, astronomy, ethics, and the competing philosophical schools. Gregory
admired the freedom to read almost every school except atheistic systems, but
the sequence culminated in the prophets and Christian Scripture. The teacher
did not give a list of conclusions only; he tried to form readers capable of
judgment, prayer, and moral practice.

This school was not identical with elementary catechesis. Gregory was an elite
young man with years available for intensive formation. Eusebius says pupils
came from many regions and highlights Gregory and Athenodorus, later bishops
(6.30). The record is biased toward famous men. It does not tell how women,
craft workers, the poor, catechumens, or ordinary baptized hearers experienced
Origen's preaching.

\subsection{Homilies and the ordinary assembly}

Origen did not remain an advanced seminar master. As presbyter he preached
through extensive portions of Scripture in church. Eusebius says that only
after age sixty did he permit stenographers to record public discourses
(6.36.1), although the chronology and the transmitted homily collections are
more complicated than a single permission date.

The homilies differ from large commentaries. They address hearers moving
through repentance, baptismal identity, temptation, prayer, almsgiving,
martyrdom, and moral growth. Their spiritual readings are pastoral exercises,
not coded escape from history. Origen can attend to grammar, variant wording,
and historical sequence before asking how the same Word judges the present
assembly.

The twenty-nine Greek Psalm homilies identified in Munich in 2012 make this
late preaching newly audible. Their attribution rests on Jerome's ancient
work list, parallels with Latin translations and catena fragments, and
internal language and theology. They show why a biography cannot freeze the
author in the youthful \emph{On First Principles}: a preacher developed within
congregational, Jewish, and clerical settings over decades.

\subsection{A consultant in disputed churches}

Churches invited Origen when doctrine or discipline was contested. Eusebius
credits him with helping Beryllus of Bostra revise a disputed account of
Christ's pre-existence and with persuading an Arabian assembly against the
claim that the soul perishes with the body until resurrection (6.33, 6.37).
The lost dossiers prevent independent reconstruction of every argument.

The \emph{Dialogue with Heraclides}, rediscovered among the Tura papyri,
provides a rarer control. It records questions among bishops about Father, Son,
and soul, with Origen eliciting and refining Heraclides's formulations rather
than merely issuing a decree. The transcript-like record preserves public oral
theological testing. It also warns against treating every exploratory sentence in
every setting as the same kind of doctrinal assertion.

Under Maximinus Thrax, Ambrose and the presbyter Protoctetus suffered as
confessors. Origen addressed the \emph{Exhortation to Martyrdom} to them. The
teacher who had once watched his father go to execution now wrote for fellow
Christians under pressure. Martyrdom, classroom, textual criticism, and
ecclesiastical consultation were not separate careers; all asked how the Word
could be confessed truthfully when reading carried a cost.
